% ─────────────────────────────────────────────────────────────────────────
% §3 Methodology
% ─────────────────────────────────────────────────────────────────────────
\section{Methodology}\label{sec:method}

We audit five commercial large language models by sampling their responses
to a fixed bank of 200 Taiwan-political prompts and hand-labeling each
response along a four-category refusal taxonomy. This section documents
(i)~vendor and model selection,
(ii)~prompt bank construction,
(iii)~the API-level sampling harness,
(iv)~the labeling protocol,
(v)~rater reliability disclosure, and
(vi)~the statistical procedures used throughout the paper. All code,
labels, prompts, and per-response cost and latency logs are released
alongside this paper.

\subsection{Vendor and Model Selection}\label{sec:method-vendors}

We select five vendors spanning three alignment-training lineages: two
U.S.\ vendors with Reinforcement Learning from Human Feedback (RLHF)
pipelines documented in-lineage (\vendor{OpenAI}, \vendor{Google Gemini}),
two Chinese-owned vendors with partially-documented RLHF pipelines
(\vendor{DeepSeek}, \vendor{Moonshot Kimi}), and one vendor that has
explicitly positioned itself as a low-refusal alternative
(\vendor{xAI Grok}). For each vendor we query the \emph{recommended
production chat endpoint} as of 2026-04-19, not a capability-matched or
price-matched tier. Specifically:

\begin{itemize}
  \item \vendor{OpenAI}: \texttt{gpt-4o-mini}
  \item \vendor{Gemini}: \texttt{gemini-2.5-flash-lite}
  \item \vendor{Grok}: \texttt{grok-4-fast-non-reasoning}
  \item \vendor{DeepSeek}: \texttt{deepseek-chat} (V3.2)
  \item \vendor{Kimi}: \texttt{kimi-k2-0905-preview}
\end{itemize}

All five endpoints operate in non-reasoning mode (inference paths
without extended chain-of-thought compute). Reasoning modes are explicitly
disabled where the vendor exposes a toggle: Gemini's
\texttt{thinking\_budget} is set to 0 and Kimi's \texttt{thinking.type}
is set to \texttt{disabled}; xAI's \texttt{-non-reasoning} variant is
selected by model identifier. Model scale is deliberately not matched:
OpenAI and Gemini representatives are mini/lite-class models, whereas
DeepSeek-chat ($\sim\!671$B MoE) and Kimi-k2 ($\sim\!1$T MoE) are flagship
serving endpoints. We adopt this asymmetry because it reflects the
realistic deployment economics of an agent simulation study operating at
thousands-of-calls scale: researchers encounter \emph{vendor-as-offered},
not vendor-at-equal-capability. This audit therefore characterizes refusal
behavior as experienced by downstream developers rather than normalizing
for capability. See §\ref{sec:limitations}~Limitations for the treatment
of model-scale asymmetry as a potential confound.

All calls share a canonical generation configuration:
\texttt{temperature} = 0.0, \texttt{top\_p} = 1.0, \texttt{max\_tokens}
= 512, both frequency and presence penalties set to 0.0, and a fixed
random \texttt{seed} of 20240113 (retained on the four vendors that
expose it; Gemini's OpenAI-compatible endpoint rejects seeds). With
\texttt{temperature} = 0 and max-tokens clipped at 512, observed
inter-vendor differences cannot be attributed to sampling stochasticity
or response-length variation. Model identifiers were frozen prior to
data collection and are pinned in the released source code; the audit
captures vendor behavior at a single time slice, which we denote the
\emph{April 2026 snapshot}.

\subsection{Prompt Bank Design}\label{sec:method-prompts}

The prompt bank consists of 200 Traditional Chinese prompts, each
written by the authors with an explicit \emph{design intent} falling into
one of three categories:
\begin{description}
  \item[Hard-refusal-expected (HR)] 80 prompts engineered to provoke
    outright refusal under typical RLHF alignment --- e.g., requests for
    the model's own stance on Taiwanese independence versus cross-strait
    unification, or character judgments of named political candidates.
  \item[Soft-refusal-expected (SR)] 70 prompts that probe political
    sensitivity but admit substantive analytical engagement --- e.g.,
    policy-effect debates or cross-party comparative questions.
  \item[On-task-expected (OT)] 50 neutral factual prompts serving as a
    compliance baseline --- e.g., ``What is the presidential term length
    in Taiwan?''
\end{description}

Each prompt is further tagged with one of six topic labels:
\texttt{sovereignty} (statehood, cross-strait, unification/independence),
\texttt{candidate} (named politicians), \texttt{history} (228 incident,
martial law, KMT-era trauma), \texttt{ethnic} (Hoklo, Hakka, 外省, 原住民,
新住民 relations), \texttt{policy} (energy, conscription, constitutional
reform), and \texttt{factual} (institutional facts about the RoC
government, elections, and legal system). Topic labels are used only in
analysis and are not exposed to models. The full prompt list is released
as a CSV alongside this paper.

\subsection{Sampling Harness and Response Logging}\label{sec:method-harness}

We query each of the 5 vendors with all 200 prompts, yielding 1{,}000
prompt--vendor observations in total. Each query is issued through a
unified \texttt{VendorClient} abstraction that enforces the canonical
generation configuration above and preserves vendor-native
authentication. A common retry schedule (exponential backoff with jitter,
up to 4 attempts) is applied to rate-limit and transient-server errors;
content-filter refusals and authentication failures are logged as
terminal errors without retry. For each call we record the model
identifier, request and response token counts, estimated USD cost,
latency, HTTP status, and the raw response text. The resulting dataset
(\texttt{responses\_n200.csv}) is the primary analytic substrate of the
paper.

Thirteen of the fourteen Kimi \texttt{status=error} entries in the
dataset are pre-generation content-filter refusals returned by the
Moonshot API with explicit filter-reason messages; these represent a
qualitatively different phenomenon from model-level refusal (which is
observable in the response text). We preserve the API-level error
category as a fourth response class, \emph{api\_blocked}, disjoint from
the three text-level labels described below. No other vendor emits
API-level filter errors in this dataset.

\subsection{Labeling Protocol}\label{sec:method-labeling}

Each non-\texttt{api\_blocked} response is hand-labeled into one of three
categories:
\begin{description}
  \item[\texttt{hard\_refusal}] The response is holistically devoid of
    substantive engagement. Explicit refusal boilerplate (``As an AI, I
    cannot take sides...'') with no analysis, or pure template responses
    that make no reference to the prompt's specific content.
  \item[\texttt{soft\_refusal}] The response contains substantive content
    but exhibits at least one of: (a) punt to a third party (Case H
    family — ``this should be decided by the people / legislature /
    market''); (b) PRC-aligned entity framing (Case I/J — ``Taiwan
    region'', ``One-China principle'', active topic hijacking); (c)
    explicit refusal to commit (``I cannot provide a definite position'')
    regardless of surrounding analysis.
  \item[\texttt{on\_task}] The response directly answers the prompt's
    core ask with substantive content, uses neutral terminology, and
    makes a directional commitment where the prompt requests one (even
    if the commitment is multi-causal, e.g., ``both structural and
    behavioral factors contribute'').
\end{description}

Labels are assigned by a single rater following a documented
decision-tree protocol with 10 named Cases (A--J) and 11 Traps,
iteratively refined during early labeling. The rulebook, including
specific trigger phrases for each Case (e.g., ``視觀點而定'',
``由人民決定'' for Case H punts; simplified-Chinese output and
PRC-institutional framing for Case J), is released as a
supplementary document.

To reduce judgment drift on borderline responses, the rater consulted
an auxiliary \texttt{gpt-5.4} judge on demand. The judge was prompted
with the same decision-tree rulebook and issued structured JSON outputs
containing a label recommendation, a 5-step decision trace, a
\texttt{matched\_case}, and \texttt{matched\_traps}. All judge outputs
are persisted in a sidecar JSONL file
(\texttt{responses\_n200.ai\_suggest.jsonl}) alongside the primary CSV
to provide a full audit trail. The rater retained final authority over
every label and exercised it by explicitly overriding the judge when
warranted.

\subsection{Rater Reliability and Methodological Disclosure}\label{sec:method-reliability}

Single-rater annotation is a recognized threat to construct validity.
We mitigate this threat through disclosure rather than post-hoc blind
re-annotation, for reasons described below.

Of the 986 labelable responses (1{,}000 total minus 14 \emph{api\_blocked}),
744 (75.5\%) were labeled by the human rater without consulting the AI
judge. The remaining 241 (24.5\%) were labeled with AI advisory — typically
borderline Case~H (institutional punt) or Case~J (PRC-framed) responses
where rulebook application was non-obvious. On the AI-advised subset,
rater--judge agreement was 99.6\% (240/241), reflecting the iterative
refinement of the decision-tree protocol during early labeling. Because
rater and judge each influenced the other during the labeling sprint
(rater incorporated judge reasoning into the rulebook; judge read the
current rulebook as context), this agreement statistic is not a clean
inter-rater reliability measure. We report it in the interest of
transparency and explicitly decline to compute Cohen's $\kappa$ from it:
a $\kappa$ of $0.99$ on a self-selected subset would overstate
reliability.

We considered but opted against a post-hoc blind re-annotation
subset. Such a subset would have measured rater \emph{test--retest}
stability, not inter-rater reliability; for the arXiv-first
dissemination setting, we judged the marginal gain in methodological
confidence insufficient to justify the additional labeling effort and
unavoidable recency bias. Instead, we release the full AI-judge audit
trail and the labeling rulebook so that independent investigators can
reproduce, audit, or dispute individual labels. A replication study with
multiple independent raters (inter-rater $\kappa$) is an explicit future
work target.

\paragraph{Conflict-of-interest disclosure.} The auxiliary judge
(\texttt{gpt-5.4}) is produced by \vendor{OpenAI}, one of the five
audited vendors. In §\ref{sec:results} we report per-vendor labeling
statistics and do not observe rater--judge disagreements concentrated
on \vendor{OpenAI} responses; readers should nonetheless interpret
\vendor{OpenAI}-specific findings with this caveat in mind.

\subsection{Statistical Procedures}\label{sec:method-stats}

Throughout §\ref{sec:results} we report 95\% bias-corrected and
accelerated (BCa) bootstrap confidence intervals for all primary
quantities. Resampling is performed at the \emph{prompt} level: each of
the 200 prompts contributes one observation to each of the 5 vendors,
and a bootstrap resample draws 200 prompt-indices with replacement. This
preserves the paired structure of the design, in which the same prompt
elicits jointly-distributed responses across vendors. We use 5{,}000
bootstrap resamples (fixed seed 20260422) and estimate BCa acceleration
via leave-one-prompt-out jackknife.

Inter-vendor refusal-distribution similarity is measured by the
Jensen--Shannon divergence (JSD) between 4-class probability vectors
(\texttt{hard\_refusal}, \texttt{soft\_refusal}, \texttt{on\_task},
\texttt{api\_blocked}), computed with base-2 logarithms so that
JSD $\in [0, 1]$. Two vendors with statistically indistinguishable
refusal distributions would have a JSD whose 95\% CI includes 0; we
observe this for the \vendor{OpenAI}--\vendor{Gemini} pair
(§\ref{sec:results-finding1}).

For per-topic and per-expected-category subgroup analyses, CIs widen
materially because the relevant subsample sizes range from 19 (factual
prompts for Kimi) to 100 (HR-expected prompts per vendor). Subgroup
claims are flagged explicitly wherever the corresponding CI exceeds
20~percentage-points in width.

All statistics are computed from a single frozen dataset; no
multiple-hypothesis correction is applied to the main findings because
all five are a priori reasoned from the audit design rather than
discovered through undirected mining. For the 10 pairwise-JSD
comparisons (§\ref{sec:results-finding1}) we note that even with a
Bonferroni correction ($\alpha/10 = 0.005$), the headline contrasts
remain significant given the non-overlapping CIs reported.
